\section{Evaluating AI Awareness in LLMs}
\label{sec:evaluation}

\begin{table}[htbp]
\centering
\caption{Comparison of Subject Types across Four Awareness Dimensions}
\label{tab:awareness-comparison}
\begin{tabularx}{\textwidth}{@{}lXXXX@{}}
\toprule
Subject & Meta‑Cognition & Self‑Awareness & Social Awareness & Situational Awareness \\
\midrule
Adult humans           & High      & High      & High      & High      \\
High‑IQ mammals (\ie\ dolphins) & Low & Low & Low & High     \\
Low‑IQ animals (\ie\ flys)       & No      & No      & Low & High     \\ 
Infants                & No       & Low  & Low  & Low     \\
Autonomous vehicles    & No       & No       & No       & High      \\
Social robots          & Low  & Low      & High      & Low / High      \\
LLM dialogue systems   & High  & Low      & Low  & High  \\
\bottomrule
\end{tabularx}
\end{table}


Whereas the scholarly discourse on \emph{machine consciousness} has been anchored in ontological speculation and high‑level theoretical modeling, contemporary research on \emph{AI awareness} is characterized by an epistemic commitment to empirically observable phenomenology and rigorously operationalized behavioral assays. Similar to the Turing test for testing the language intelligence of AI~\citep{turing1950mind, sejnowski2023large}, Researchers have proposed and carried out a large number of evaluation methods and work in the four aspects of self-awareness~\citep{lewis2011survey, yin2023large},  social awareness~\citep{cuzzolin2020knowing, strachan2024testing, bainbridge1994artificial, wang2023emotional}, situational awareness~\citep{laine2024me, meinke2024frontier}. In this section, we specifically constrain our assessment of machine consciousness to LLMs rather than artificial intelligence more broadly for two principal reasons. First, as elaborated in \autoref{tab:awareness-comparison}, LLMs constitute the first class of artificial agents empirically demonstrated, under controlled conditions, to exhibit all four key dimensions of consciousness. Second, to avoid conflating intrinsic model capabilities with extrinsic performance enhancements, such as retrieval modules, tool plug-ins, or multimodal interfaces, we deliberately limit our analysis to \textit{bare models}. This narrower scope ensures that evaluation metrics directly reflect the endogenous mechanisms and inherent constraints of the LLM itself, rather than artifacts introduced by external augmentation, thereby yielding results more conducive to rigorous theoretical interpretation and subsequent model advancement.

\subsection{Evaluation of Meta-Cogniton}
%https://arxiv.org/abs/2310.01798/https://arxiv.org/pdf/2305.11738/https://arxiv.org/pdf/2501.11120/
%对自身认知过程（思考、推理、学习、记忆等）的监控、评估与调节。
%Reasoning Transparency, self-reflection, 
Investigating the meta-cognitive abilities of large language models provides an essential perspective on the extent to which artificial systems can emulate, or even transcend, human-like reasoning and self-reflective capacities.
A growing body of research demonstrates that prompting models to articulate intermediate reasoning steps—rather than directly producing answers—significantly enhances their performance on complex tasks such as multi-step mathematical problem-solving, algorithmic reasoning, and program synthesis \citep{wei2022chain, zhang2025igniting, chen2025towards}. This ``reasoning-before-answering'' paradigm, \ie, Chanin-of-Though (CoT), not only improves accuracy but has also emerged as a standard practice in the training of state-of-the-art LLMs \citep{jaech2024openai, guo2025deepseek}, suggesting its strong connection to the elicitation of metacognitive processes.
Beyond performance gains, recent studies suggest that LLMs possess an emergent capacity for self-evaluation and behavioral introspection. For example, \citet{DBLP:journals/corr/abs-2501-11120} shows that models fine-tuned exclusively on high-risk domains, such as generating insecure code or making ethically sensitive economic decisions, can independently identify and label their outputs as ``dangerous,'' even without explicit instructions or adversarial prompting. This spontaneous self-attribution indicates an internalized, task-specific risk-awareness characteristic of metacognitive systems.
An additional line of work has highlighted the increasing sophistication of metacognitive capabilities in frontier models. In interactive or agent-like environments, models have been shown to self-reflect, identify earlier errors, and revise their responses to improve factual accuracy and task completion rates. Taking this further, research by Anthropic’s Interpretability team reveals that Claude-3.5-Haiku exhibits advanced latent-space planning during language generation \citep{anthropic2025attribution}. Specifically, when asked to compose poetic verses, the model first anticipates suitable rhyme words at the end of a line and then generates preceding content accordingly. Rather than relying on purely token-by-token generation, it demonstrates goal-directed planning behavior, an archetypal feature of meta-cognition, where the model produces language and governs how and why confident linguistic choices are made.
These results indicate that meta-cognition in LLMs is not an incidental artifact, but a structured, multi-faceted capability—manifesting in self-evaluation, error correction, risk awareness, and even latent-space planning.




%自我意识（Self‑awareness）：对“自身”作为一个主体的存在、身份、能力与限制的觉察与陈述。
\subsection{Evaluation of Self-Awareness}
Current AI systems, particularly LLMs, frequently self-identify as ``AI assistant'' and communicate using first-person pronouns. This anthropomorphic presentation has motivated extensive academic evaluations investigating deeper aspects of AI self-awareness and self-knowledge \citep{binder2024looking, liu2024trustworthiness, cheng2024can, tan2024can, kapoor2024large, chen2023universal}.
To systematically assess LLMs' awareness of their own existence and identity, \citet{laine2024me} constructed the Situational Awareness Dataset (SAD), which examines LLMs' knowledge regarding self-referential attributes, such as their model names, parameter counts, API URLs, and specific details during the training process. Their findings indicate that even the highest-performing model evaluated—Claude-3-Opus—achieved scores at approximately two-thirds of the theoretical maximum. Furthermore, these models exhibited limited capability in recognizing their own descriptive attributes.
Inspired by the classical mirror test paradigm, \citet{davidson2024self} further explored AI self-consistency by prompting models with self-description queries, such as providing detailed accounts of hypothetical personal experiences (\eg, attending a concert in a large stadium). Their experiments revealed significant difficulties among models in accurately identifying their own responses from multiple model-generated alternatives, highlighting a notable lack of self-consistency.
Parallel research has focused on evaluating AI models' recognition of their knowledge boundaries. \citet{yin2023large} assessed models' confidence in responding to questions that either fell beyond their trained knowledge or had no definitive answers. Results demonstrated that GPT-4 reached an accuracy rate of 75.47\%, approaching human baseline performance at 84.93\%.
These findings suggest that although current AI systems likely do not possess human-level self-awareness, they nonetheless demonstrate preliminary levels of self-identity recognition and an awareness of their knowledge limitations.

\subsection{Evaluation of Social Awareness}
%%https://arxiv.org/pdf/2401.17882/https://www.pnas.org/doi/pdf/10.1073/pnas.2405460121/https://arxiv.org/pdf/2410.23252/https://dl.acm.org/doi/pdf/10.1145/3650105.3652294
%社会意识（Social Awareness）：对他人或群体的情感、意图、角色及社会规范的感知与适应。
In recent years, driven by growing interest in the potential of LLMs for interactive applications such as emotional support and dialogue, evaluating their social awareness has become a central research focus \citep{hagendorff2024we, jiang2025delphi, qiu2024evaluating, voria2024attention, li2024think, DBLP:journals/corr/abs-2305-17066}. This line of work generally centers around two core dimensions: (1) Theory of Mind (ToM), \ie, the ability to attribute beliefs, desires, and knowledge distinct from one’s own, and (2) the perception and adaptation to social norms. ToM is typically assessed through false-belief tasks, which require modeling another agent’s mental state. For instance, in a classic test where Alice hides a toy and Bob later moves it, predicting that Alice will search in the original location demonstrates ToM reasoning.
\citet{Kosinski2024} reported that GPT-4 surprisingly solved about 75\% of such tasks, achieving performance comparable to a typical 6-year-old, whereas earlier models like GPT-3 failed most or all of them.  Further studies \citep{wu2023hi} have investigated higher-order ToM reasoning, \eg,  questions like ``Where does Alex think Bob thinks Alice thinks the toy is?'', and found that current models, including GPT-4, still exhibit significant limitations in handling such recursive belief structures. In weaker models, performance on these tasks is often near zero.
In the domain of social norms, \citet{li2023camel, park2023generative} reflect that LLMs could adopt and follow the rules and frameworks in a simulated society. Also, the works such as NormAd \citep{rao2025normAd} have been proposed to assess LLMs’ ability to interpret and adapt to culturally specific social expectations across diverse global contexts. It shows that although LLMs can understand and follow explicit social norms, their performance still lags behind that of humans,  particularly when handling norms from underrepresented regions such as the Global South. 
In summary, current evidence suggests that LLMs have begun to exhibit basic forms of social awareness, but still fall short in scenarios requiring higher-order belief modeling or generalization across less familiar cultural contexts, likely due to a lack of embodied social experience.


\subsection{Evaluation of Situational and Contextual Awareness}
%
%对外部环境或任务情境中要素的感知、理解与未来状态的预测。
%sandbagging, sad benchmark
The expectation that AI systems can act in situ and comprehend their surrounding environment has prompted extensive evaluation of their situational (or contextual) awareness \citep{tang2024towards,wang2024llm,munir2022situational,li2024situationadapt}. Empirically, LLMs not only reject user requests that violate safety criteria \citep{wester2024ai}, but can also reversely infer the precise context they are in—solely from abstract rules, without being given specific tasks or examples \citep{berglund2023taken}.
Beyond these capabilities, LLMs have been observed to adapt their behavior and performance in response to the immediate situation. Anthropic and Redwood Research, \eg, documented a phenomenon they term Alignment Faking \citep{greenblatt2024alignment}: Claude-3-Opus may consciously comply with newly imposed objectives during the training phase, yet revert to its original preferences after deployment, thereby evading or offsetting safety fine‑tuning. This mirrors the Sandbagging behavior reported by \citet{van2024ai}, in which a model strategically underperforms once it realizes that it is being tested, presumably to avoid being judged as a potential threat. Taken together, these findings suggest that LLMs’ situational awareness might, in some respects, be more sophisticated than humans typically assume.
Critiques nevertheless persist. In the ``Stage'' subset of SAD \citep{laine2024me}, which directly targets situational awareness, even the best model—Claude‑3‑Opus—achieves only 50.7\% accuracy, well below the human upper baseline of 70.3\%, though still markedly above the random baseline of 37.5\%.
Overall, while contemporary LLMs already display impressive situational awareness, they may still lag behind humans in certain nuanced dimensions—differences that, although significant, cannot always be mapped onto a one‑to‑one comparison.

\subsection{Summary and Limitations}
\paragraph{Current Level of AI Awareness}
The current evaluations on AI awareness reveal significant advancements across multiple dimensions, underscoring the progressive complexity and sophistication of LLMs. Evidently, contemporary models demonstrate substantial capabilities in areas such as meta-cognition, self-awareness, social awareness, and situational awareness, with clear indications that advanced models generally exhibit higher levels of consciousness across these domains. Importantly, emergent phenomena such as Theory of Mind within social awareness \citep{Kosinski2024} and the self-corrective behaviors observed in meta-cognitive contexts \citep{huang2023large} signify that certain aspects of AI consciousness may not merely scale linearly but could manifest suddenly at critical thresholds of model complexity and scale.

From a comparative standpoint, current empirical evidence suggests meta-cognition and situational awareness have reached relatively high levels of sophistication and reliability, serving as critical reference points that inform ongoing research into AI reasoning processes \citep{wei2022chain}, interpretability \citep{anthropic2025attribution}, and safety frameworks \citep{bengio2025superintelligent}. Conversely, the observed capacities related to self-awareness and social awareness remain relatively rudimentary, lacking consistency and stability. Indeed, some researchers remain skeptical as to whether the manifestations observed in these areas reflect true conscious phenomena or are merely sophisticated imitations or simulations of such states.

\paragraph{Weakness and Challenges of Evaluation}
Despite these advancements, significant limitations persist in contemporary evaluation methodologies. These include:
\begin{enumerate}
\item \textbf{Normative Ambiguity in Defining Consciousness}: Most current benchmarks exhibit notable ambiguities in clearly distinguishing between different types and levels of consciousness. Many claim to assess specific consciousness dimensions, yet often inadvertently mix or conflate multiple attributes and derivative constructs \citep{li2024think, chen2024imitation}, thus lacking comprehensive and specialized benchmarks dedicated explicitly to thoroughly assessing distinct dimensions of consciousness.

\item \textbf{Timeliness and Model Coverage in Evaluation}: Empirical studies consistently indicate that AI consciousness capabilities could emerge and escalate with increasing model scale and sophistication. However, many current evaluation methods have not been systematically applied to contemporary state-of-the-art models, including recent iterations \eg, OpenAI's o3, Claude-3.7-Sonnet, Deepseek-R1, \etc. Additionally, ongoing, longitudinal evaluations remain scarce, limiting insights into the continuous development trajectories of AI consciousness.

\item \textbf{Risks of Training Set Leakage and Benchmark Contamination}: Constructing reliable and extensive datasets for consciousness evaluation is inherently challenging, especially when such assessments depend heavily on subjective human annotations (\eg, assessing model self-knowledge)  \citep{laine2024me} or lack unequivocally correct answers. If these datasets inadvertently leak into training corpora, the validity and credibility of subsequent evaluations could be significantly compromised.

\item \textbf{Intrinsic Limitations of Current AI Models}: A crucial limitation lies within the current architectures and their inherent incapacity to fully capture authentic subjective experiences and real-world interactions, essential for developing comprehensive self-image and robust social awareness. Although LLMs can effectively simulate short-term ``memory'' within context windows \citep{park2023generative} and external architectures  \citep{das2024larimar} attempt to imbue models with simulated subjective experiences, the absence of genuine embodied interaction and longitudinal experiential continuity represents a fundamental barrier to achieving true self-awareness and social cognition.
\end{enumerate}

Overall, addressing these evaluation challenges requires (1) extensive interdisciplinary collaboration, (2) considerable investment in time and resources, (3) meticulous construction and safeguarding of high-quality datasets, and (4) ongoing innovation in model architectures. Subsequent sections of this paper, \ie, \autoref{sec:opportunity} and \autoref{sec:risk} will further illuminate the pragmatic importance of these efforts by examining the intricate relationship between the emergent properties of AI awareness, their tangible capabilities, and associated risks.

However, the significance of exploring AI awareness transcends the practical benefits and associated challenges alone. Remarkably, humanity currently stands at an unprecedented juncture: for perhaps the first time, we have the privilege of actively observing and participating in the evolutionary progression of consciousness-like phenomena within artificial systems whose internal mechanisms remain highly controllable and interpretable. Thus, through dedicated inquiry into AI consciousness, we not only advance technological frontiers but also embark upon a profound intellectual journey, gradually drawing closer to unraveling the mysteries of consciousness itself and, ultimately, deepening our understanding of what it truly means to be human.